\section{秦：商鞅入秦与一场高代价的重组}

秦孝公求贤时，关中的道路已经不只是通往西方的边路。渭河两岸的耕地、雍城以来的旧都邑和向东越过函谷的出口，使秦拥有可供经营的腹地；但它仍须面对一个不利的事实：中原的魏先一步把旧晋的资源组织成强国，秦在河西和东部通道上屡受挤压。前362年的\eventanchor{WQ-0362-WEIQIN}{战国·少梁之战}正是这种压力的可见一刻。战争并未自动推出一套改革方案，却让秦的统治者更难继续依靠旧的封邑关系和临时征发来同魏、齐、楚等国竞争。

商鞅不是从秦国贵族网络里长出来的人。传统材料称他出自卫国公族旁支，曾在魏相公叔痤门下任事；《史记》叙述魏未能任用他，随后他西入秦廷。无论这段离魏入秦的细节有多少经后世剪裁，跨国流动本身并不奇特：战国的游士、将领和说客在诸国间携带的，不只是言辞，还有对户籍、军队、官吏和邻国强弱的观察。秦孝公选择任用这样一位外来者，也意味着他愿意以触犯本国既得关系的代价，换取一次重新分配功劳、土地和控制权的机会。

前356年的\eventref{WQ-0356-SHANGYANG}{战国·商鞅第一次变法}，在传世叙事中因而成为这场选择的起点。它把什伍互保与连坐、以军功授爵、奖励耕战、编户与对游食的限制等措施连在一起。它们的共同方向并不是一句抽象的“富国强兵”，而是让一个家庭的户口、生产、治安责任和出战所得更容易被国家登记、核验并赏罚。对依靠世卿和宗族关系取得地位的人，这意味着原有的权利要被军功和法令重新估价；对普通农户，则意味着取得田宅与爵位的路径可能更清楚，同时也更深地进入赋役、军役和相互监督的网络。

\begin{textwindow}[变法者的辩辞]
“治世不一道，便国不法古。”

出处为《史记·商君列传》。司马迁将这句话放在商鞅初说秦孝公、为改变旧制辩护的场合；它凝练了后世所理解的变法逻辑，却出自汉代传记对战国论辩的重构，不能视为前356年朝廷辩论的逐字记录。
\end{textwindow}

这种重组不能只归结为一位改革者的意志。孝公提供了授权和政治庇护，秦既有的关中农业、旧有行政基础和边地军事经验提供了可用资源；魏、楚、韩等国已经在各自处境中尝试军政整顿，则构成了持续的外部压力与可比较的经验。商鞅的方案之所以能在秦推进，正因为这些条件在同一时刻相互咬合；它所激起的贵族抵触，也说明法令的执行从来不是纸面上的单向命令。

前350年，秦自雍迁都咸阳，并在传统所称的第二次改革中继续调整制度。这里的\eventanchor{WQ-0350-QINREFORM}{战国·秦迁咸阳与第二次变法}不宜被读成一次突然完成的“现代化”：迁都把政治中心移到更接近东进通道和渭河交通的位置；合并乡邑为县、设置令丞，以及开阡陌封疆、校正斗桶权衡丈尺等叙述，则指向一个更细密的行政、田制与市场尺度。它们若能落实，国家便不必只在出兵时才接触地方，而能在田界、度量、县署和道路的日常往来中持续取得信息和资源。

这正是改革的高代价所在。县制和法令压缩了世袭贵族留在地方的裁量空间；军功爵把风险、奖赏和服从更紧地捆在一起；以户为单位的登记和连坐，也把治理失灵的风险部分转嫁给相邻家庭。秦得到的是更可调用的人口、粮食与兵员，却不是一部无摩擦的机器。制度要靠地方官吏、基层家庭、军功受益者和统治集团反复执行，任何一环的抵制、规避或滥用都可能改变它的效果。

\subsection{变法前的关中：继承、边线与东部都城}

商鞅到来前的秦并非一张白纸。前419年前后，秦灵公已在河西同魏氏争渡口、城邑，\eventanchor{WQ-0419-QIN-LING-BORDER}{战国·秦灵公经营河西边线}；战役名、胜负和每座城的归属没有完整年表，却足以看出关中向东的压力早已存在。灵公卒后，秦简公即位，\eventanchor{WQ-0415-QIN-JIANGONG}{战国·秦简公即位}，其后迁治雍地、整理关中城邑的传统，见于\eventanchor{WQ-0415-QIN-JIAN-ADMIN}{战国·秦简公整顿关中城邑}。迁治的动机与工程不可坐实，但都城离边线的远近，本就是军政如何看待河西竞争的一部分。

继承使这条边线一度更难处理。前400年简公卒、惠公继位，\eventanchor{WQ-0400-QIN-JIAN}{战国·秦简公卒、惠公继位}；前387年惠公卒而出子立，\eventanchor{WQ-0387-QIN-HUI}{战国·秦惠公卒、出子继位}。灵公身后与出子继承相关的争议在前388年前后渐成焦点，\eventanchor{WQ-0388-QIN-LING-DEATH}{战国·秦灵公身后继承纷争}；至前385年出子被废、献公复位，\eventanchor{WQ-0385-QIN-XIAN}{战国·秦献公复位}。政变的支持者和清理范围没有完整记录，故只能谨慎说献公重组旧臣与军政，\eventanchor{WQ-0385-QIN-XIANGONG-CLEANUP}{战国·秦献公清理旧臣}，而不能把它写成一份逐人的肃清名单。

传世材料还把献公时废止从死列为制度转折，\eventanchor{WQ-0385-QIN-END-HUMAN-SACRIFICE}{战国·秦献公废止从死}。禁令的覆盖范围、延续时间均无同时代文书可核，它至多表明王权干预旧俗的做法被后人郑重记住。比这更能触摸到的，是献公在栎阳经营东部支点，\eventanchor{WQ-0384-QIN-XIANGONG-CAPITAL}{战国·秦献公经营栎阳}；城址考古可帮助确认其位置，却不能替我们复原每日政务。前378年前后的东向试探，\eventanchor{WQ-0378-QIN-XIANGONG-EAST}{战国·秦献公东向反攻}，尚未扭转魏强秦弱，却让关中统治集团认识到战线、都城与征发不可分开。献公晚年留下的河西压力与栎阳军政，构成\eventanchor{WQ-0368-QIN-XIANGONG-LEGACY}{战国·献公遗留的东进压力}；少梁前夕秦魏反复交涉、备战，\eventanchor{WQ-0366-QIN-WEI-NEGOTIATION}{战国·秦魏围绕河西交涉}，说明战争并非改革故事的一个方便开场，而是已经逼近的条件。

\subsection{求贤令不是一张空白诏书}

前364年石门之战中，秦传称斩首六万，\eventanchor{WQ-0364-SHIMEN}{战国·石门之战}；整齐的首级数不可化作人口统计，却显示秦已把对魏反攻当作可用军功语言。前361年孝公即位、求贤，\eventanchor{WQ-0361-XIAOGONG}{战国·秦孝公即位求贤}，给了改变政策的政治入口。流传下来的求贤令文字出自传世叙事，不能当作一纸原诏；它所引起的三晋、卫、宋士人向栎阳流动，仍可作为\eventanchor{WQ-0361-QIN-SEEK-TALENT-RESPONSE}{战国·秦孝公求贤引士人西来}来理解。景监推荐公孙鞅，\eventanchor{WQ-0361-JINGJIAN-RECOMMEND}{战国·景监荐公孙鞅}；卫人公孙鞅由魏入秦，\eventanchor{WQ-0360-SHANGYANG-MIGRATION}{战国·公孙鞅自魏入秦}。荐举、入秦路线和初见言辞都经过列传铺陈，但外来者借求贤而进入秦廷这一政治事实，不必依赖那些对话的逐字真实性。

约前359年，商鞅与孝公议变法，\eventanchor{WQ-0359-SHANGYANG-COUNSEL}{战国·商鞅与孝公议变法}。他提出什么次序、孝公经历几次试探，不能由《商君列传》的成篇辩辞逐句还原；比较可靠的是，改革必须先取得君主持续授权，才会碰到宗室、太子与地方。什伍与户籍连坐的推进，\eventanchor{WQ-0358-QIN-HOUSEHOLD-REGISTRATION}{战国·什伍户籍与连坐的推进}，以及告奸、匿奸赏罚，\eventanchor{WQ-0356-QIN-INFORMING-LAW}{战国·告奸匿奸赏罚}，把看见治安与看见户口绑在一起。条文由《商君书》和列传归纳，未必为同日颁布的完整法典；它们所指的方向，却是令基层家庭彼此承担可被核验的责任。

军功重估宗室与贵族的位置，\eventanchor{WQ-0356-QIN-CLAN-NOBLES}{战国·军功爵限制旧贵族特权}。太子驷触法、其师受刑的故事，\eventanchor{WQ-0355-QIN-HEIR-PUNISHMENT}{战国·太子周边受刑}，刑罚名称和日期略有异文，却清楚地保存了改革如何把太子集团卷入冲突。它不是一则“法在君上”的抽象寓言：改革若要被相信不是只约束下层，便必须触到本来最难触动的人；它也因此为孝公死后的人事清算积下了怨隙。

\subsection{县、田与兵：制度并不在同一天完成}

前350年以后，县邑、军功爵与田宅授予相互扣合，\eventanchor{WQ-0351-QIN-COUNTY-REGRANT}{战国·县邑军功与田宅授予}。睡虎地简牍属于较晚时期，只能旁证成熟的文书实践，不能把它倒灌为前351年全秦县份的统一表格；但它提醒我们，县级官吏以律令、户籍和传舍协调田赋、兵役，\eventanchor{WQ-0348-QIN-COUNTY-OFFICIALS}{战国·秦县级行政的文书运作}，正是在这些细小程序里增加国家对地方的可见度，也同时扩大吏员的裁量。

兵符是另一种可见的凭证。约前四世纪中叶的杜虎符等器物，\eventanchor{WQ-0350-QIN-TIGER-TALLY}{战国·杜虎符所见调兵凭证}，能证明调兵需要可验证的符信，不能确定某一件器物在何次战事中使用。它与咸阳迁都、田界和县署共同说明：改革追求的不是一纸宏大宣言，而是让命令经过物件、官吏和道路仍可被核对。

前340年商鞅诱俘魏将公子卬、取河西要地，\eventanchor{WQ-0340-SHANGYANG-WEI}{战国·商鞅取河西}。战场细部和诱敌过程带有列传修辞，河西胜利却使商鞅受封商、於一带，\eventanchor{WQ-0340-SHANGYANG-FENGSHANG}{战国·商鞅受封商於}。私人封邑将改革者的权力、军功与新法的声望拴得更紧：它一面奖励了可见的军事成果，一面也使此人日后难以同新君分开。

\subsection{改革者死后，制度仍在向外走}

前338年孝公卒，惠文君继位，\eventanchor{WQ-0338-HUIWEN}{战国·秦惠文君继位}；商鞅失势，\eventanchor{WQ-0338-SHANGYANG-DEATH}{战国·商鞅失势被杀}。他逃往封地、聚众抵抗，\eventanchor{WQ-0338-SHANGYANG-FLIGHT}{战国·商鞅逃往封地}，终于在彤地失败、被处死并及其族，\eventanchor{WQ-0338-SHANGYANG-TONGYI-BATTLE}{战国·商鞅彤地兵败}。旅店拒宿、逃亡路线和司法程序皆有传奇化的成分；值得注意的不是替故事补足细节，而是个人被清除后，县制、军功等路径并未一概废去，\eventanchor{WQ-0339-QIN-NOBLE-OPPOSITION}{战国·商鞅死后新法延续}。新君保留什么、改动什么都不应说得过满，但制度已同军政利益相连，确比一位客卿的命运更能延续。

惠文王随后以宗室与客卿并用。前337年前后，樗里疾等进入相将决策，\eventanchor{WQ-0337-QIN-HUIWEN-APPOINTMENTS}{战国·惠文王调整相将}；前333年阴晋改名宁秦，\eventanchor{WQ-0333-NINGQIN}{战国·阴晋改名宁秦}；前329年又渡河取汾阴、皮氏一带，\eventanchor{WQ-0329-FENYIN}{战国·秦取汾阴皮氏}。改名意图、城邑地望与占领时间均有考辨，不能把名称本身写成已稳固的边界；它们仍显示秦开始把河西胜利接入东向的连续压力。樗里疾等向魏河东据点施压，\eventanchor{WQ-0327-CHULIJI-WEI-CAMPAIGN}{战国·樗里疾攻魏河东}，至前325年秦君称王，\eventanchor{WQ-0325-QIN-KING}{战国·秦惠文君称王}，军事、官署与王号已在相互支撑一种更主动的国家姿态。

扩张也检验制度的承载范围。前328年前后，樗里疾、甘茂等进入军政核心，\eventanchor{WQ-0328-CHULIJI}{战国·樗里疾甘茂入军政核心}；秦取巴蜀后设郡守、张若等治理新占地，\eventanchor{WQ-0315-SHU-ADMIN}{战国·秦在蜀设置郡守}，并长期维护入蜀栈道与山地交通，\eventanchor{WQ-0315-SHU-ROAD}{战国·秦整治蜀道}。建置先后、道路工程和官员职掌不全可定，不能把后世蜀道的全貌放回这一年。前301年蜀侯煇反秦、司马错平定，\eventanchor{WQ-0301-SHU-HUI-REBELLION}{战国·蜀侯煇反秦}，反倒提醒人们：兼并和设官不是一劳永逸，新地区仍会以地方权力和交通成本考验关中的命令。

秦武王时初置左右丞相，\eventanchor{WQ-0309-QIN-TWO-CHANCELLORS}{战国·秦置左右丞相后分工协调}。左右相实际如何分权不可等同后世宰相制度，制度意图却很清楚：宗室、客卿和军队必须有人协调。前306年昭襄王诛庶长壮后，宣太后与魏冉集团巩固摄政，\eventanchor{WQ-0306-QIN-REGENCY-PURGE}{战国·昭襄王朝的摄政整合}；罪名与涉案范围记载简略，不能以此复原一场完整政变。前300年前后，秦又在蜀道沿线设置转输节点、维护山路，\eventanchor{WQ-0300-QIN-SHU-ROUTE-ADMIN}{战国·秦经营蜀道转输节点}。亭障名称和年份不宜虚拟，山地道路需要反复投入劳役这一点却无需神化。商鞅改革因此不是终点；它留下的是此后每一位统治者都得继续维修、检验并可能加重使用的行政与动员路径。

\begin{sourcenote}[从商鞅故事回到材料边界]
\sourcebook{史记·秦本纪}与《商君列传》写得最完整，也最富戏剧性；《魏世家》《樗里子甘茂列传》《司马错列传》补出河西、相将与蜀地的国别线索。它们成书远在事件身后，保存改革次序与政治冲突的强大传统，却不能让每一段游说、每一项条文和每一个动机都成为可直接复原的现场记录。《商君书》保存农战、赏罚和国家控制的思想语言，但篇章的年代、作者和编纂过程存在争议，不能把全书逐条当作商鞅在前356年或前350年颁行的法令。睡虎地、里耶等秦简与杜虎符，让人看见较晚法律文书、县级行政、考核和调兵凭证如何运转；栎阳城址、蜀道与巴蜀地方史材料则帮助检验空间条件。它们证明的是这些制度传统后来能长期展开，不能倒推每一项成熟做法都出自商鞅个人、也不能精确判定其初行日期。
\end{sourcenote}

从这层意义上说，商鞅改革是后来秦国家能力的一项重要条件，而不是秦统一的单一原因。它帮助形成了能够持续征发、记录、赏罚和补给的制度路径；后来的统治者必须继续维护、修订并把它用于新的土地和战争，关中资源、地理位置、对手之间难以协调的联盟，以及长期战争中的将领和官吏，同样决定了结果。到\eventref{WQ-0293-YIQUE}{战国·伊阙之战}、\eventref{WQ-0260-CHANGPING}{战国·长平之战}以及\eventref{QIN-0230-UNIFICATION}{秦·统一六国}时，这种能力才在不同条件下显示出累积效力；它既不能把所有胜利追溯为商鞅一人的设计，也不能抹去改革曾经改变国家与社会相互关系的分量。

\begin{turningpoint}[制度得失：能力不是一个人的遗产]
这场重组缓解了秦在强邻竞争中难以稳定征兵、征粮和控制地方的难题，代价则由旧贵族、地方社会与被更严密编入国家的家庭共同承担。它把可动员资源和行政信息连接得更紧，也把高强度征发与严厉赏罚留作后续统治者可以使用、也可能过度使用的工具。秦的上升应放在这条制度路径、关中条件与多国竞争的合力中理解，而不应写成单一人物决定历史的直线故事。
\end{turningpoint}
